\section{Discussion}
To test the functionality and endurance of the final system, it was run for 2 hours straight, while a human stood by to record and log errors as well as intervene to fix said errors.
Ideally the test would have run for significantly longer, but due to the limited number of available robots, the test duration was limited to 2 hours to avoid blocking other project groups for too long from continuing their own work.

% Cycles, runtime, and errors
The 2 hour test completed a total of 81 cycles (see Appendix~\ref{app:2HourResults} Table \ref{tab:2hour-test-results}) with a total of 31 errors, where each cycle consists of the robot putting a cube into storage, and returning a cube from a randomly chosen slot to the input area (see Whole system test in Appendix~\ref{videos}). Thus an average cycle time of $~1.48min$, with an error detected every $~2.61$ cycles in average.
% Total errors
This revealed the greatest error producing subsystems to be Vision and Storage with $45\%$ and $48\%$ of total errors respectively. This is equivalent to an error occurring in these subsystems in $17\%$ and $19\%$ of cycles during the 2-hour runtime.

% Storage errors
Within the errors attributed to Storage, inability to locate slot was the most frequently occurring with $23\%$ of all errors being attributed to this error type. Meaning this occurs in $9\%$ of cycles on average.
It is very likely that some of these errors can be explained by light from external sources (e.g. the Sun) interfering with LDR readings. The possibility of adding a shield for the LDR to block the external interference and improve performance, should be investigated further.

% Vision errors
Although Storage had the greatest percentage of errors, the most common error type was camera being unable to read data from QR code, which took up $29\%$ of all errors, while the third most common was camera taking more time than usual to read the data, which took up $16\%$ of all errors. Both of these error types belongs to the Vision subsystem. These error types thus occurs in $11\%$ and $6\%$ of cycles on average respectively. % The most frequent errors are therefore related to reading data from QR codes, making it an obvious area for future improvement.
Adding some form of cover to protect the camera from external light, while also adding controlled lighting, should make QR-code reading less likely to fail. This would make the images taken by the camera more consistent and reduce the risk of having too much or too little light to read the QR code.
Alternatively, since the subsystem errors were exclusively related to reading the encoded data, compromising the ability to read the QR codes, potentially marking them as unknown, in favor of only obtaining position data, would remove these errors and thus improve system performance.

% Can be removed if necessary
%In regards to robot communication, alternative methods of more low level control could have been implemented to extract further data perform and more customized movement, if given the time. However this is not a necessity for the system to work and operate successfully, and would thus be a much lower priority, when compared to other issues with the system.









%\subsection{Perspective}